\section{\rev{Numerical Verification and Sensitivity}}\label{sec:completed-verification}
\begin{revision}
The numerical assessment examines the discretization, treatment of curved boundaries, reconstruction of interface fields, and parameter sensitivity. A manufactured solution provides an exact reference for the regular operator, while an independent finite-element discretization is used to assess the curved-interface solution. Synthetic images with known geometries provide a separate evaluation of the measurement procedure.
\end{revision}

\subsection{Regular manufactured Poisson problem}
The source-derived quartic test uses $\phi=r^4+z^4+r^2z^2$ and the axisymmetric source $18r^2+16z^2$. On grids $17\times19$, $33\times37$, $65\times73$, and $129\times145$, recorded RMS errors were $1.82155502\times10^{-3}$, $4.44529637\times10^{-4}$, $1.09730635\times10^{-4}$, and $2.72549081\times10^{-5}$. Consecutive observed orders were 2.03482, 2.01831, and 2.00938.
\begin{revision}
Discretization error is evaluated as the RMS difference between the numerical and exact potentials at interior nodes. For halved mesh spacings, the observed order is $q=\log(e_h/e_{h/2})/\log 2$. As shown in Figure~\ref{fig:regular-check}, the FD error decreases systematically with refinement, and the observed order approaches two. An independently assembled bilinear Q1 finite-element discretization exhibits the same asymptotic behavior. These results are consistent with second-order accuracy for the regular axisymmetric Poisson operator.
\revisionfigure{0.98}{01_regular_poisson_fem.png}{Regular-grid convergence. Errors are measured against the exact manufactured potential; $h$ is the geometric mean of the radial and axial mesh spacings in model length units.}{fig:regular-check}

\subsection{Curved-boundary refinement}
The effect of curved-boundary representation is examined by comparing immersed and staircase discretizations on grids of $31\times45$, $45\times65$, $61\times89$, and $87\times129$. To maintain consistent spatial sampling, the potential is evaluated at the same 56 gas-side locations on all grids. The change between successive levels is defined by
\[
D_\ell=\left[\sum_{j=1}^{56}\bigl(\phi_\ell(\mathbf{x}_j)-\phi_{\ell-1}(\mathbf{x}_j)\bigr)^2\right]^{1/2}.
\]
Figure~\ref{fig:curved-check} shows that $D_\ell$ decreases from 10.334 to 4.264 and 2.549~V for the immersed discretization. The corresponding staircase values are 25.731, 7.107, and 11.267~V. Apparent orders calculated from these changes are 2.854 and 1.429 for the immersed treatment, compared with 4.148 and $-1.280$ for the staircase treatment. The immersed formulation therefore exhibits a consistent decrease in inter-level variation over the tested grids, whereas the staircase representation produces a reversal at the finest pair of levels. This comparison characterizes the refinement sensitivity of the curved-boundary solution.
\revisionfigure{0.98}{02_curved_boundary_refinement.png}{Curved-boundary refinement. Potential changes use fixed gas-side points; apparent orders are calculated from consecutive changes.}{fig:curved-check}

\subsection{Independent fixed-cone finite-element comparison}
An independent comparison is performed between the immersed finite-difference (FD) solver and a body-fitted, piecewise-linear P1 finite-element (FEM) discretization of the same charge-free electrostatic problem. The domain radius and axial extent are 12.5 and 10~mm, respectively. The rounded cone has a half-angle of $49.29^\circ$, an apex at $z=8.6$~mm, and a cap radius of $50~\mu\mathrm{m}$. A potential of 1000~V is prescribed on the cone, with the outer boundary grounded. Three FD grids, $65\times97$, $129\times193$, and $257\times385$, are compared with successively refined FEM meshes. Both solutions are evaluated at common gas-side locations and at 41 interface positions over $z=3$--6~mm.

The finest-level relative RMS differences are 0.024\% for potential, 0.226\% for the vector electric field, 0.381\% for the interface-normal field, and 0.817\% for Maxwell pressure $p_E=\epszero E_n^2/2$. The projected-voltage difference is 1.071\%, and the projected-residual difference is 2.518\% when normalized by the RMS of the centered capillary pressure. Identical curvature, arclength weights, pressure projection, and voltage-amplitude definitions are used in the two calculations. The agreement in Figure~\ref{fig:fixed-reference} consequently extends from the primary potential solution to the derived interface stresses and projected residuals.
\revisionfigure{0.98}{03_fixed_cone_fd_fem.png}{Fixed-cone FD--FEM comparison. (a) Gas-point differences. (b) Interface-normal field. (c) Maxwell pressure at 1000~V. (d) Amplitude-projected residual.}{fig:fixed-reference}

\subsection{Angle and parameter sensitivity}
The sensitivity of the angle diagnostics is assessed over candidate half-angles of $46^\circ$--$52^\circ$ by varying mesh resolution, domain size, cap size, flank-sampling window, and outer-boundary condition. Figure~\ref{fig:angle-sensitivity} compares the projected-residual landscape with the exponent obtained by fitting $E_n^2\propto\rho^p$, where $\rho$ denotes distance from the virtual cone vertex. Under the analytical Taylor far-field condition, the residual-minimizing angle varies from $48.0^\circ$ to $50.0^\circ$, whereas the $p=-1$ crossing remains within $49.07^\circ$--$49.25^\circ$. Replacing this boundary condition with a grounded outer boundary increases the residual and shifts the exponent curve. The two diagnostics therefore resolve different features: the crossing measures consistency with the imposed Taylor scaling, while the residual minimum reflects the finite-domain stress balance.
\revisionfigure{0.98}{04_sensitivity_landscapes.png}{Angle sensitivity. Residual minima and field-exponent crossings respond differently to numerical settings.}{fig:angle-sensitivity}

Initialization sensitivity is further examined using five angle--apex-radius starting pairs for Powell optimization on the $31\times45$ grid. The final half-angles span $40.199^\circ$--$42.791^\circ$, and the apex-radius parameters range from 0.0704 to 0.0887~m. Despite these distinct endpoints, the projected RMS residuals remain within 0.01257--0.01333~Pa (Figure~\ref{fig:multistart}). This weak variation in objective value indicates an angle--radius ambiguity in the fitted parameterization. Accordingly, the residual magnitude and the recovered geometry provide complementary information on the optimization result.
\revisionfigure{0.98}{05_optimizer_multistart.png}{Multistart optimization. Specified starts and final parameter pairs are paired with their final projected residuals.}{fig:multistart}

\subsection{Synthetic-image measurement test}\label{sec:synthetic-completed}
The accuracy of image-based geometry extraction is assessed on analytic cone silhouettes. Each silhouette is a conical profile with a spherical-cap apex, drawn at $0.02$~mm/pixel for a needle of radius $0.4$~mm, and is pushed through the preprocessing, edge-detection, contour-extraction and flank-fitting chain that the imaging toolkit provides for recorded frames. The test matrix crosses three half-angles ($45^\circ$, $49.29^\circ$ and $55^\circ$) with four Gaussian noise levels whose standard deviations are 1.3, 4.0, 12.7 and 40.0 greyscale levels, corresponding to nominal signal-to-noise ratios of 40, 30, 20 and 10~dB.

Image scale and axis location are prescribed as calibration data: the symmetry axis is taken from the needle centreline, which is mounted on the optical axis and therefore known a priori, and the noise realization uses random seed 20260922. The absolute angle error is defined as $|\alpha_{\mathrm{extracted}}-\alpha_{\mathrm{true}}|$, and the profile nRMSE is the RMS radial difference from the analytic profile normalized by the visible cone length. Over the two lowest noise levels, at which the acceptance criteria are asserted, the largest angle error is $0.027^\circ$ and the largest profile nRMSE is 0.62\%, and all twelve cases satisfy the prescribed $1^\circ$ and 2\% criteria (Figure~\ref{fig:synthetic-check}). Heavy noise degrades the profile but not the fitted slope: at the highest noise level the $55^\circ$ silhouette returns a profile nRMSE of 10.5\% while its angle error remains $0.010^\circ$.

These values quantify the implementation accuracy of the extraction chain, not the uncertainty of a measurement. Because every silhouette is rendered from the same analytic profile that the fit is compared against, the test exercises the software path --- rendering, thresholding, contour tracking and least-squares flank fitting --- and does not represent illumination variation, contact-line pinning, optical distortion, or the finite resolution of a recorded frame. The extraction chain is therefore qualified for the numerical round trip only, and the comparisons of Section~\ref{sec:experimental-procedure} remain subject to their own scope statement.
\revisionfigure{0.98}{06_synthetic_holdout.png}{Synthetic-image accuracy. Each marker is one analytic-image case; dashed lines show the preset $1^\circ$ and 2\% acceptance limits. Both axes are logarithmic.}{fig:synthetic-check}
\end{revision}
